% A document only XeTeX can set.
%
% `fontspec` refuses to load under pdfTeX -- it raises a fatal package error
% and stops -- so this document is how the corpus shows a pdfTeX-only
% distribution refusing cleanly rather than producing a wrong page. The
% Unicode below is the second reason: pdfTeX has no way to set Greek and
% Cyrillic from a Unicode source without a package that maps them.
%
% No font is named. There are no system fonts in a worker, so a document that
% asked for one by name would fail on the browser distributions for a reason
% that has nothing to do with the engine, and this example is about the
% engine.
\documentclass[11pt]{article}

\usepackage{fontspec}
\usepackage{amsmath}

\title{A Document That Needs XeTeX}
\author{The Komodoc Corpus}
\date{}

\begin{document}
\maketitle

\section{Unicode, straight from the source}

Greek, in the source as Greek rather than as \verb|\alpha|: αβγδε ζηθικ λμνξο.

Cyrillic, likewise: Привет, мир.

Mathematical Greek still comes from the maths fonts, and is unaffected:
\begin{equation}
  \alpha^{2} + \beta^{2} = \gamma^{2}.
\end{equation}

\section{Why this one is here}

A distribution whose only engine is pdfTeX must refuse this document with the
error \verb|fontspec| gives, and that refusal must reach the reader as a
diagnostic rather than as a hang or a blank page. That is the whole of what
this example tests.

\end{document}
